\chapter{Courses and rooms}
\label{ch:courses}

\section{Create a course}

\begin{center}
\cmd{!course create <code> "<title>" [--semester HS26] [--template <name>]}
\end{center}

\begin{description}[leftmargin=7em,style=sameline,font=\normalfont\ttfamily]
  \item[<code>] A short, unique code such as \cmd{CS101}. Letters, digits, \cmd{.}, \cmd{_} and
        \cmd{-}, at most 32 characters, no spaces. The bot stores it in capitals, and you can
        type it in any case afterwards.
  \item["<title>"] The course title. Put it in quotes if it contains spaces.
  \item[--semester] \cmd{HS} (autumn semester) or \cmd{FS} (spring semester) plus the
        two-digit year, e.g.\ \cmd{HS26}. Without it, the current semester is used
        (August--January: HS, February--July: FS).
  \item[--template] Use a saved layout instead of the default one (\autoref{ch:templates}).
\end{description}

\begin{you}
!course create CS101 "Introduction to Computer Science" --semester HS26
\end{you}
\begin{botreply}
Course \textbf{CS101} is ready: https://matrix.to/\#/!abc...\\
Rooms: Announcements, General, Questions, Staff\\
You are invited to the space and all rooms; accept the invitations in Element.\\[4pt]
\textbf{Next steps} \dots
\end{botreply}

The space is called \enquote{\texttt{CS101 Introduction to Computer Science (HS26)}}. You
become the course \emph{owner}: the bot gives you full rights in every room, and only you
(and bot admins) can change the course's structure.

\subsection*{The default rooms}

\begin{center}
\begin{tabularx}{\textwidth}{@{}lLL@{}}
\toprule
\textbf{Room} & \textbf{Who can join} & \textbf{Who can write} \\
\midrule
Announcements & everyone in the course space & staff only (owner and TAs) \\
General & everyone in the course space & everyone \\
Questions & everyone in the course space & everyone \\
Staff & only the owner and TAs (by invitation) & owner and TAs \\
\bottomrule
\end{tabularx}
\end{center}

All rooms are end-to-end encrypted. New members can read the history from before they
joined. Students join a room by opening the course space in Element and clicking the room.

\subsection*{Permissions inside the rooms}

Matrix rooms have \emph{power levels}. The bot sets them for you:

\begin{center}
\begin{tabular}{@{}lll@{}}
\toprule
\textbf{Who} & \textbf{Power level} & \textbf{Shown in Element as} \\
\midrule
Course owner & 100 & Admin \\
TA & 50 & Moderator \\
The admin bot & 100 & Admin \\
Students & 0 & Default \\
\bottomrule
\end{tabular}
\end{center}

Moderators can remove messages, kick people and invite. Only admins change room settings.

\begin{warning}[Don't remove the bot]
The bot has to stay in the course space and its rooms with its admin rights. Without them it
can't invite students, create groups, archive the course or add rooms. Don't kick the bot or
lower its power level.
\end{warning}

\section{See your courses}

\begin{itemize}
  \item \cmd{!course list}: every course you own or are a TA of, with semester and role.
        Archived courses are marked. Bot admins see all courses.
  \item \cmd{!course info CS101}: title, semester, owner, TAs, the number of students (joined
        or invited), the rooms, how many signup links nobody has used yet, and a link to the
        space.
\end{itemize}

\section{Teaching assistants}
\label{sec:tas}

\begin{you}
!course ta add CS101 @ta.one:matrix.dmi.unibas.ch
\end{you}

The TA must have an account already. Write the full Matrix ID; in Element you see it by
clicking a person's name. The bot invites the TA to every room of the course (including Staff
and the group rooms) and makes them moderator there.

\cmd{!course ta remove CS101 @ta.one:matrix.dmi.unibas.ch} takes the moderator rights away and
removes the person from staff-only rooms. They stay in the normal rooms like any member.

\begin{center}
\begin{tabularx}{\textwidth}{@{}Lcc@{}}
\toprule
\textbf{In a course \dots} & \textbf{TA} & \textbf{Owner} \\
\midrule
see course info, signup links of the course & \faCheck & \faCheck \\
invite people, create signup links for the course & \faCheck & \faCheck \\
roster sync, also with \cmd{--remove} & \faCheck & \faCheck \\
create, fill and post to exercise groups & \faCheck & \faCheck \\
appoint or remove TAs & -- & \faCheck \\
add or remove rooms & -- & \faCheck \\
export the layout as a template & -- & \faCheck \\
archive, roll over & -- & \faCheck \\
\bottomrule
\end{tabularx}
\end{center}

\section{Rooms}

\begin{you}
!room add CS101 "Lab sessions" --topic "Questions about the weekly lab"
\end{you}

A new room is open to everyone in the course space, like General. With \cmd{--private} only
staff are invited and nobody else can join, like the Staff room. The topic is optional.

\cmd{!room remove CS101 "Lab sessions"} takes the room out of the course space. The room
itself still exists with its history and members; it just no longer appears in the course.

\section{The course over a semester}
\label{sec:lifecycle}

\begin{center}
\begin{tikzpicture}[
  phase/.style={draw=accent,fill=accentlight,rounded corners=3pt,font=\footnotesize\sffamily,
    align=center,minimum height=1.1cm,text width=2.6cm},
  arr/.style={-{Stealth[length=2.5mm]},accent,thick}]
  \node[phase] (c) {\textbf{create}\\\cmd{!course create}};
  \node[phase,right=0.3cm of c] (t) {\textbf{staff}\\\cmd{!course ta add}};
  \node[phase,right=0.3cm of t] (i) {\textbf{students}\\\cmd{!invite}};
  \node[phase,right=0.3cm of i] (g) {\textbf{groups}\\\cmd{!groups ...}};
  \node[phase,below=0.9cm of g] (r) {\textbf{next semester}\\\cmd{!course rollover}};
  \node[phase,left=0.3cm of r] (a) {\textbf{close}\\\cmd{!course archive}};
  \draw[arr] (c) -- (t); \draw[arr] (t) -- (i); \draw[arr] (i) -- (g); \draw[arr] (g) -- (r);
  \draw[arr] (r) -- (a);
\end{tikzpicture}
\end{center}

After a rollover, the new course starts again at \enquote{students}: invite the new list.

\subsection*{Rollover: the same course next semester}

\begin{you}
!course rollover CS101 FS27
\end{you}

\begin{enumerate}
  \item The current course is renamed to \cmd{CS101-HS26}, so the code \cmd{CS101} is free
        again. Its rooms, messages and members stay as they are.
  \item A new course \cmd{CS101} for FS27 is created with the same rooms and settings, the same
        owner and the same TAs. Students are not carried over.
\end{enumerate}

Afterwards invite the new students with \cmd{!invite CS101}, and archive the old course.

\subsection*{Archive: close a course}

\begin{you}
!course archive CS101-HS26
\end{you}

After you confirm the plan:

\begin{itemize}
  \item students can no longer write anywhere in the course; staff still can,
  \item every room gets the prefix \cmd{[Archiv HS26]},
  \item the space is moved into the server's archive space, if the admins set one up,
  \item with \cmd{--kick-students}, students are also removed from all rooms.
\end{itemize}

Everything stays readable for the members. Use \cmd{--dry-run} to see the plan without
creating one.

\begin{note}[Deleting a course]
The bot doesn't delete courses: chat history is often needed later (questions about grades,
plagiarism cases). If a course was created by mistake, ask a bot admin.
\end{note}
